% !TeX program = lualatex

% Warning: template parts using for loops give an error when they are indented.
%See when having trouble using graphicx with invoice: http://tex.stackexchange.com/questions/61113/using-graphicx-with-invoice-results-in-strange-pdf-output

\documentclass[paper=a4,fontsize=11pt,DIV=12]{scrlttr2}
\usepackage[T1]{fontenc}
\usepackage{lmodern}
\usepackage[gen]{eurosym}
\usepackage{ucs}
\usepackage[utf8x]{inputenc}
\usepackage{microtype}
\usepackage{dcolumn}
\usepackage{booktabs}
\usepackage[english]{babel}
\usepackage{underscore}
%\usepackage{float, caption}
\usepackage{graphicx}
%\usepackage{invoice_bammerlaan}
\usepackage{fp}
\usepackage{changepage} % Allow moving table slightly to left
\usepackage{fontspec} % Load fontspec for Roboto Light font
\setmainfont[BoldFont={Roboto-Regular.ttf}]{Roboto-Light.ttf}
\usepackage{mathastext}
\usepackage{wallpaper}
\usepackage{etoolbox}

\LoadLetterOption{DINmtext}
\KOMAoptions{enlargefirstpage=true,
            fromalign=right,
            fromphone=true,
            fromemail=true,
            fromurl=true,
            fromlogo=true,
            backaddress=false,
            parskip=half*,
            firstfoot=false,
            foldmarks=off}

%%%%%%%%%%%%%%%%%%%%%%%%%%%%%%%%%%%%%%%%%%%%%%%%%
% Set global variables here:

% (Company) Name
% \setkomavar{title}{Invoice}
\usepackage{datatool}
\DTLsetseparator{ = }
\DTLloaddb[noheader, keys={key,val}]{komavarvalues}{komavarvalues.dat}
\newcommand{\komavarfromfile}[1]{\DTLfetch{komavarvalues}{key}{#1}{val}}


\setkomavar{fromname}{\komavarfromfile{fromname}}
% Address
\setkomavar{fromaddress}{\komavarfromfile{fromaddress}}

% Phone number
\setkomavar{fromphone}{\komavarfromfile{fromphone}}
% E-mail address
\setkomavar{fromemail}{\komavarfromfile{fromemail}}
% Website
\setkomavar{fromurl}{\komavarfromfile{fromurl}}

%%%%%%%% Files %%%%%%%%

% Path to LaTeX template (so graphics can also be found when run from another directory) (REPLACE WITH YOUR OWN):
\graphicspath{ {/home/user/yourdirectoryhere} }

% Path to background image
% \ULCornerWallPaper{1}{achtergrond_facuur.pdf}

% Path to logo
\setkomavar{fromlogo}{\includegraphics[height=3cm]{invoice_main_logo}}


%%%%%%%%%%%%%%%%%%%%%%%%%%%%%%%%%%%%%%%%%%%%%%%%%



\setkomavar{firstfoot}{%
	\parbox[t]{\textwidth}{\footnotesize
		\begin{tabular}[t]{l@{}}%
			\multicolumn{1}{@{}l@{}}{Chamber of commerce: \KVKNummer{}}
		\end{tabular}%
		\hfill
		\begin{tabular}[t]{l@{}}%
			\multicolumn{1}{@{}l@{}}{IBAN: \usekomavar{frombank}}
		\end{tabular}%
		\hfill
		\begin{tabular}[t]{l@{}}%
			\multicolumn{1}{@{}l@{}}{VAR number: 
				\BTWNummer{}}
		\end{tabular}%
	}%
}


% Uncomment for European date entries. Also uncomment line starting with "\parsedate}" below. --- see: http://tex.stackexchange.com/questions/54594/tex-capacity-exceeded-while-parsing-a-date-string
%\def\parsedate#1{\edef\temp{#1}%
%	\expandafter\parsedateX\temp\relax}
%\def\parsedateX #1-#2-#3\relax{%
%	\def\dueyear{#1}%
%	\def\duemonth{#2}%
%	\def\dueday{#3}}

% Prevent indentation of signature: 
\renewcommand*{\raggedsignature}{\raggedright}

\providebool{hasdiscount}


\begin{document}

    \begin{letter}{To:\\
			\textbf{@{owner['name']}}\\
			@{owner['full_name']}\\
%+ for a in owner['address']:
@{a}\\
%-
		}
		
		
%		\setkomavar{subject}{%
%			@{notes}}
		\setkomavar{subject}{Invoice}		
		\setkomavar{invoice}[\textbf{Invoice number}]{@{id}}
		
		\opening{Dear @{owner['full_name'].split(' ')[0]},}
		
        Please see the breakdown of this invoice below and contact me if you have any questions. 
        
        If paying by check, please make it out to \usekomavar{fromname}.
        
        Finally, please pay within 30 days of receipt of this invoice.

		% @{terms['desc']}



%+ for e in entries:
\ifbool{hasdiscount}
{}
{ \setbool{hasdiscount}{@{str(e['discount'] is not None).lower()}} }
%-

        %% TODO: Add optional Discount Column using above boolean!!
            %% Don't forget the header, table entries, and the footer calculations!

\ifbool{hasdiscount}{
  \newcommand{\discountcolhead}{Discount}
  \newcommand{\discountcoldata}[1]{#1}
  \newcommand{\colspec}{
                l % Date
				p{15em} % Description
				D{,}{,}{2} % Amount
				l % Unit price
				l % Amount without tax
                l % Pre-tax Discount
				l % Tax rate
				l % Tax amount
				l % Total price with taxes
  }
}{
  \newcommand{\discountcolhead}{}
  \newcommand{\discountcoldata}[1]{}
  \newcommand{\colspec}{
                l % Date
				p{13em} % Description
				D{,}{,}{2} % Amount
				l % Unit price
				l % Amount without tax
                @{} % Hidden
				% l % Tax rate
				l % Tax amount
				l % Total price with taxes
  }
}

\renewcommand{\arraystretch}{1.4}

	\begin{adjustwidth}{-4em}{} % Move table slightly left	
		\expandafter\tabular\expandafter{\colspec}
			\multicolumn{1}{p{2.5em}}{\flushleft Date} & % Date
			\multicolumn{1}{p{14em}}{\flushleft Description} & % Description
			\multicolumn{1}{p{3em}}{\flushleft Qty} & % Amount
			\multicolumn{1}{p{3em}}{\flushleft Unit price} & % Unit price
			\multicolumn{1}{p{4em}}{\flushleft Subtotal} & % Amount without tax
            \ifbool{hasdiscount}{
                \multicolumn{1}{p{4em}}{\flushleft \discountcolhead} & % Pre-tax Discount
                % I think have to do this weird boolean finagling here due to the use of multicolumn ¯\_(ツ)_/¯
            }
            {
            }
			% \multicolumn{1}{p{1.5em}}{\flushleft Sales\\Tax} & % Sales Taxed?
			\multicolumn{1}{p{2em}}{\flushleft Tax\\amount} & % Tax amount
			\multicolumn{1}{p{4em}}{\flushleft Total}\\ \midrule % Total price with taxes
%+ for e in entries:
@{ e['date'].strftime('%m/%d') } & % Date
% Comment the above and uncomment the below for European dates:
%\parsedate{@{e['date']}} \dueday-\duemonth-\dueyear & % Date
@{e['description'] if '(' not in e['description'] else "\\newline(".join(e['description'].split("("))} & % Description

% This line may have weird syntax highlighting, since it is trying to abide by both Python's syntax
% and whatever processing rules are in gcinvoice... 
% I tried to figure out a better way, such as f strings, but that fails due to the right bracket
% being detected as the end of the embedded python expression itself. ¯\_(ツ)_/¯
@{"%.2f" % float(e['qty'])} & % Number of Units

\$ @{e['price']} & % Unit price
\$ @{_currencyformatting(float(e['amount_net']) + float(e['amount_discount']),precision=2)} & % Amount without tax and before discount

% Pre-tax Discount
\ifbool{hasdiscount}
{ \$ ( @{e['amount_discount']} ) & }
{}

% @{'Y' if 'taxtable' in e else 'N'} & % Tax rate
@{'\$ ' + str(e['amount_taxes']) if 'taxtable' in e else '-'} &
%@{owner['full_name']} 
\$ @{e['amount_gross']} \\ % Total price with taxes
%-
			\midrule
            &	Sub-total	& &	 & \$ @{_currencyformatting(float(amount_net)/0.85,precision=2)}	& \$ (@{_currencyformatting(float(amount_net)/0.85 * 0.15,precision=2)} )       &  \$ @{amount_taxes}     &	\$ @{amount_gross}\\ \cmidrule{8-8}
%			&	+ BTW &	&  &     &   &    &	\$ @{amount_taxes} \\ \cmidrule{8-8}
			&	\large\textbf{Amount due} &		&  &     &    &   &	\textbf{\scalebox{1.2}{\$ @{amount_gross}}} \\

% If discount:			
			
			% &	\multicolumn{2}{l}{Discount} &  & &   &   &	\$ @{cformat(Decimal('0.1') * amount_net_)} \\ \cmidrule{8-8}
			% &	\textbf{Total minus discount} &  & &   &   &	& \$ @{cformat(amount_gross_ - Decimal('0.1') * amount_net_)}
		\endtabular
	\end{adjustwidth}				
		\closing{Yours sincerely,}
	\end{letter}
\end{document}